\documentclass[11pt]{article}
\usepackage[margin=1in]{geometry}
\usepackage{amsmath}
\usepackage{booktabs}
\usepackage{longtable}
\usepackage{array}
\usepackage{hyperref}
\usepackage{url}

\title{Memo 12: Spatial Distribution Analysis for FVCOM-MP Cases D and E\\
       FSK Settling and Effective Grain-Size Coupling, Days 10 and 30}
\author{FVCOM-MP WaterPACT Investigation}
\date{May 2026}

\begin{document}
\sloppy
\setlength{\tabcolsep}{4pt}
\maketitle

\section*{Purpose}

This memo records a new spatial-distribution comparison using the day-10 and
day-30 \texttt{OUTPUT\_d} and \texttt{OUTPUT\_e} MAT files.  The paired cases
use the same cohesive FSK sediment setup and the same small MP size class.  The
primary question is whether turning on the MP-side equivalent floc or grain
size treatment changes the MP fields.

\section*{Files and Command}

The analysis inputs were:
\begin{itemize}
  \item \texttt{FVCOM\_MP\_test\_run/OUTPUT\_d/spatial\_distribution\_v2\_d\_day10.mat}
  \item \texttt{FVCOM\_MP\_test\_run/OUTPUT\_e/spatial\_distribution\_v2\_e\_day10.mat}
  \item \texttt{FVCOM\_MP\_test\_run/OUTPUT\_d/spatial\_distribution\_v2\_d\_day30.mat}
  \item \texttt{FVCOM\_MP\_test\_run/OUTPUT\_e/spatial\_distribution\_v2\_e\_day30.mat}
\end{itemize}

The analysis was run with:
\begin{verbatim}
python FVCOM_MP_test_run/PYTHON/analyze_spatial_distribution_v2.py ^
  --days 10 30 --cases d e --diff-hi d --diff-lo e ^
  --output-base FVCOM_MP_test_run/PYTHON/output/spatial_distribution_v2_de
\end{verbatim}

The script was updated so the difference pair is configurable.  For this memo,
difference maps and \texttt{diff\_summary\_stats.csv} use
\texttt{case d - case e}.  The outputs were written to:
\begin{itemize}
  \item \texttt{FVCOM\_MP\_test\_run/PYTHON/output/spatial\_distribution\_v2\_de/day10}
  \item \texttt{FVCOM\_MP\_test\_run/PYTHON/output/spatial\_distribution\_v2\_de/day30}
\end{itemize}
The run generated 108 PNG maps across the two day folders.

\section*{Extraction Window}

All four MAT files contain 24 output records.  The day-10 files span
simulation day 9.0000 to 9.9570, and the day-30 files span simulation day
29.0000 to 29.9570.  These are final-24-hour averages before the target day
endpoints, matching the v2 workflow used in Memo 09.

\section*{Paired Setup}

Cases \texttt{d} and \texttt{e} are a cleaner pair than the Memo 09
\texttt{c}--\texttt{d} comparison because the active sediment setup is the
same.  Both cases use cohesive sediment, \texttt{SED\_WSET\_METHOD=FSK\_CONC},
\texttt{SED\_SD50=0.010} mm, and \texttt{SED\_WSET=0.086} mm s$^{-1}$.
The plastic setup is also the same except for the effective-floc switch:

\begin{table}[h]
\centering
\caption{Active setup difference between cases \texttt{d} and \texttt{e}.}
\begin{tabular}{lll}
\toprule
Item & Case \texttt{d} & Case \texttt{e} \\
\midrule
Sediment file & \texttt{generic\_sediment\_d.inp} & \texttt{generic\_sediment\_e.inp} \\
Sediment type & cohesive & cohesive \\
Settling method & \texttt{FSK\_CONC} & \texttt{FSK\_CONC} \\
\texttt{SED\_SD50} & 0.010 mm & 0.010 mm \\
\texttt{SED\_WSET} & 0.086 mm s$^{-1}$ & 0.086 mm s$^{-1}$ \\
Plastic file & \texttt{generic\_plastic\_d.inp} & \texttt{generic\_plastic\_e.inp} \\
\texttt{PLAST\_FLOC\_EFFECTIVE} & \texttt{T} & \texttt{F} \\
MP axes & 0.010, 0.002, 0.002 mm & 0.010, 0.002, 0.002 mm \\
\bottomrule
\end{tabular}
\end{table}

In the aggregation closure, the MP incorporation threshold uses the minimum MP
dimension.  For this setup, the relevant MP size is therefore 0.002 mm, or
2 $\mu$m.  The volume-equivalent MP diameter is approximately 3.42 $\mu$m.

\section*{Key Spatial Statistics}

Table~\ref{tab:stats} summarizes depth-averaged domain means for the fields
most relevant to the FSK/effective-size question.  Case \texttt{e} is the
matched control with \texttt{PLAST\_FLOC\_EFFECTIVE=F}; case \texttt{d} has
the effective-size coupling turned on.

\small
\begin{longtable}{llrrr}
\caption{Depth-averaged domain means from the generated summary CSV files.}
\label{tab:stats}\\
\toprule
Day & Field & Case \texttt{e} & Case \texttt{d} & D/E or change \\
\midrule
\endfirsthead
\toprule
Day & Field & Case \texttt{e} & Case \texttt{d} & D/E or change \\
\midrule
\endhead
\bottomrule
\endfoot
10 & Sediment & $9.555{\times}10^{-1}$ & $9.555{\times}10^{-1}$ & no change \\
10 & $W_{\mathrm{floc}}$ (m s$^{-1}$) & $5.014{\times}10^{-5}$ & $1.533{\times}10^{-4}$ & $3.06{\times}$ \\
10 & $d_{\mathrm{floc,eff}}$ ($\mu$m) & $16.01$ & $26.41$ & $1.65{\times}$ \\
10 & $\lambda_a$ (s$^{-1}$) & $5.887{\times}10^{-1}$ & $3.181{\times}10^{-1}$ & $0.54{\times}$ \\
10 & $\lambda_d$ (s$^{-1}$) & $1.394{\times}10^{-1}$ & $3.559{\times}10^{-2}$ & $0.26{\times}$ \\
10 & MP total (kg m$^{-3}$) & $5.499{\times}10^{-7}$ & $6.037{\times}10^{-7}$ & $1.10{\times}$ \\
10 & MP aggregated (kg m$^{-3}$) & $2.721{\times}10^{-7}$ & $5.307{\times}10^{-7}$ & $1.95{\times}$ \\
10 & MP dispersed (kg m$^{-3}$) & $2.778{\times}10^{-7}$ & $7.302{\times}10^{-8}$ & $0.26{\times}$ \\
10 & Aggregated MP fraction & $0.494$ & $0.855$ & $+0.361$ \\
10 & Bottom MP mass (kg m$^{-2}$) & $4.965{\times}10^{-11}$ & $7.886{\times}10^{-10}$ & $15.9{\times}$ \\
\midrule
30 & Sediment & $9.703{\times}10^{-1}$ & $9.703{\times}10^{-1}$ & no change \\
30 & $W_{\mathrm{floc}}$ (m s$^{-1}$) & $8.105{\times}10^{-5}$ & $2.078{\times}10^{-4}$ & $2.56{\times}$ \\
30 & $d_{\mathrm{floc,eff}}$ ($\mu$m) & $26.12$ & $39.05$ & $1.50{\times}$ \\
30 & $\lambda_a$ (s$^{-1}$) & $5.057{\times}10^{-1}$ & $2.705{\times}10^{-1}$ & $0.53{\times}$ \\
30 & $\lambda_d$ (s$^{-1}$) & $1.199{\times}10^{-1}$ & $2.746{\times}10^{-2}$ & $0.23{\times}$ \\
30 & MP total (kg m$^{-3}$) & $1.041{\times}10^{-6}$ & $1.315{\times}10^{-6}$ & $1.26{\times}$ \\
30 & MP aggregated (kg m$^{-3}$) & $5.143{\times}10^{-7}$ & $1.140{\times}10^{-6}$ & $2.22{\times}$ \\
30 & MP dispersed (kg m$^{-3}$) & $5.264{\times}10^{-7}$ & $1.758{\times}10^{-7}$ & $0.33{\times}$ \\
30 & Aggregated MP fraction & $0.494$ & $0.863$ & $+0.369$ \\
30 & Bottom MP mass (kg m$^{-2}$) & $9.516{\times}10^{-11}$ & $1.369{\times}10^{-9}$ & $14.4{\times}$ \\
\end{longtable}
\normalsize

\section*{Main Findings}

\subsection*{The sediment fields are unchanged}

The \texttt{d-e} difference statistics show exactly zero difference in
\texttt{sed\_davg}, \texttt{sed\_surf}, and \texttt{sed\_bot} for both days.
This is important: the analysis isolates the MP-side effective-size switch
rather than a sediment-regime change.  It also means that the MP response below
is not caused by a different sediment concentration field.

\subsection*{The FSK equivalent size is active and visible}

With \texttt{PLAST\_FLOC\_EFFECTIVE=T}, case \texttt{d} uses the active FSK
settling signal to compute the floc size used by the MP interaction path.  The
depth-averaged diagnostic settling speed is about 3.06 times case \texttt{e}
on day 10 and 2.56 times case \texttt{e} on day 30.  The corresponding
domain-mean diagnostic Stokes-equivalent floc size increases from 16.0 to
26.4 $\mu$m on day 10 and from 26.1 to 39.1 $\mu$m on day 30.

Case \texttt{e} still has a diagnostic \texttt{d\_floc\_eff} field in the MAT
files because the extractor back-calculates it from the saved settling
velocity.  However, with \texttt{PLAST\_FLOC\_EFFECTIVE=F}, the model's
aggregation kernel uses the static sediment size \texttt{SED\_SD50=0.010} mm
inside the MP-floc interaction.

\subsection*{The small MP class remains below the incorporation threshold}

The aggregation code uses the minimum MP dimension for the Wu-style
incorporation threshold, so the relevant MP size is 2 $\mu$m.  For a 10 $\mu$m
static sediment size, the threshold
\[
  d_{\mathrm{MP,max}} =
  -0.0002 d_{\mathrm{floc}}^2 + 0.36 d_{\mathrm{floc}}
\]
is about 3.58 $\mu$m.  For the case-\texttt{d} mean effective floc sizes of
26--39 $\mu$m, the threshold is about 9--14 $\mu$m.  The 2 $\mu$m MP class is
therefore allowed to aggregate in both cases, but the effective-size case has a
much larger margin.  This is different from Memo 09, where the 100 $\mu$m MP
class was too large for the small effective flocs.

\subsection*{Aggregation rate decreases, but aggregated MP fraction increases}

Turning on the effective-size treatment reduces the domain-mean aggregation
rate by about 46\% on both days.  It also reduces the disaggregation rate more
strongly: by about 74\% on day 10 and 77\% on day 30.  The lower
disaggregation rate dominates the state variables: the aggregated MP fraction
rises from about 0.494 in case \texttt{e} to 0.855--0.863 in case
\texttt{d}.  The aggregated MP concentration approximately doubles, while the
dispersed MP concentration drops to roughly one-quarter to one-third of the
control.

\subsection*{Bottom MP mass is highly sensitive to the switch}

The bottom MP mass is much larger when the effective-size treatment is on:
about 15.9 times the control on day 10 and 14.4 times the control on day 30.
The bottom aggregated MP mass remains zero in the generated summaries, so this
increase appears in the bottom dispersed MP mass diagnostic.  The result still
indicates a strong transport and deposition response associated with the
effective-size treatment.

\section*{Interpretation}

The answer to the main question is yes: turning on the equivalent grain or floc
size treatment has a clear impact in this paired test.  Because cases
\texttt{d} and \texttt{e} have identical sediment concentration fields, the
impact is attributable to the MP-side use of the FSK-derived settling and
equivalent-size information.

The mechanism is not a simple increase in every aggregation metric.  The
effective-size treatment makes the MP interaction use larger effective flocs
and a stronger settling signal.  For the very small MP class used here, the MP
remains comfortably below the floc incorporation threshold, so aggregation is
not shut off.  Instead, the rate diagnostics show lower instantaneous
aggregation but much lower disaggregation, which shifts the modeled MP state
toward aggregated suspended mass and much larger bottom mass.

This paired \texttt{d}--\texttt{e} comparison is therefore a useful production
test for the new FSK/effective-size pathway.  It shows that the new pathway is
active, spatially resolved, and dynamically important for small MPs.

\section*{Recommended Next Checks}

\begin{enumerate}
  \item Inspect the day-30 difference maps for \texttt{lambda\_a\_davg},
        \texttt{lambda\_d\_davg}, \texttt{agg\_fraction\_davg},
        \texttt{mp1\_agg\_davg}, and \texttt{bot\_mass}; these are the fields
        where the switch has the clearest spatial footprint.
  \item Add a direct model output diagnostic for the floc size actually used by
        the aggregation kernel and for $d_{\mathrm{MP,max}}$.  The current MAT
        files include a useful post-processed equivalent size, but the
        aggregation-kernel size is still inferred from the switch logic.
  \item Repeat the same paired analysis for day 90 once
        \texttt{spatial\_distribution\_v2\_d/e\_day90.mat} files are
        available.
\end{enumerate}

\end{document}
